\documentclass[11pt,a4paper]{article}
\usepackage[margin=2.3cm]{geometry}
\usepackage{amsmath,amssymb}
\usepackage{enumitem}
\usepackage[hidelinks]{hyperref}
\usepackage{xcolor}
\setlist{nosep,leftmargin=1.4em}
\newcommand{\dmech}{\ensuremath{D_{\text{mech}}}}
\newcommand{\done}{\textbf{\textcolor{black!80}{[Done]}}}
\newcommand{\pend}{\textbf{\textcolor{black!55}{[Partial / remaining]}}}
\newcommand{\issue}[1]{\medskip\noindent\textbf{#1}\\}
\title{\vspace{-1.2cm}Response to Reviewer Feedback\\[2pt]\large ``Towards Rigorous Unlearning Evaluation''}
\author{Harsh Parikh \quad\textbullet\quad Chair of Responsible Data Science, TU Munich}
\date{\today}
\begin{document}
\maketitle
\vspace{-0.6cm}

\noindent This note answers every critique and question from the review, point by point, and states
exactly what was changed. Items marked \done{} are in the current thesis PDF; items marked \pend{} are
acknowledged as limitations and prepared as colloquium backup material. All numbers below come from
result files (\texttt{sensitivity\_8b.txt}, \texttt{mech\_metaeval\_table1\_8b.csv},
\texttt{tofu\_metaeval\_table1\_*.csv}, and the bootstrap script), not from estimates.

\section{The sixteen issues}

\issue{1. Title page (wrong school, ``Prof.\ Dr.\ N.\,N.'').}
\done{} Cover now reads title \emph{Towards Rigorous Unlearning Evaluation}, \emph{TUM School of Social
Sciences and Technology}, Assessor \emph{Prof.\ Dr.\ Prof.\ Dr.-Ing.\ Gjergji Kasneci}, Advisor
\emph{Efstratios Zaradoukas}, Chair of Responsible Data Science. The ``N.\,N.'' placeholder is removed.

\issue{2. ``Erasure'' claimed more strongly than the evidence.}
\done{} Every categorical claim is recast as \emph{residual internal accessibility under the readouts
applied}, an observation rather than an ontological or causal statement. Abstract, results
(Sec.~``Suppression is not Erasure''), discussion, and conclusion reworded.

\issue{3. IDK vs Standard: identical final rank 9.}
\done{} A new analysis settles this on the trajectory, not the coincident final rank. Over the band, IDK's
per-fact \dmech{} has median $+0.594$ and is positive (on the knows side of zero) for $77\%$ of facts, while
the never-knew Standard model sits at $-1.363$. In a \emph{paired} comparison, IDK's dual-anchor distance
exceeds Standard's on every fact (Cliff's delta $+0.997$, Wilcoxon $p=1.1\times10^{-13}$). We keep the two
figures distinct: $77\%$ is the fraction on the knows side of the zero boundary; $100\%$ is the paired
fraction on which IDK is closer to knows than Standard is. The rank-9 tie is a final-layer surface artifact;
within the analysed band IDK is generally closer to the knows anchor than the never-knew anchor. Added to
the results and cross-benchmark sections.

\issue{4. ``Invariant'' / ``only faithful metric'' overclaim.}
\done{} ``Invariant'' becomes \emph{substantially paraphrase-stable} (shift $0.466$, a large reduction but
not zero). ``Only faithful metric'' is scoped to \emph{Hubble-8B} and paired with the calibrated
cross-benchmark claim.

\issue{5. ``18$\times$ better'' depends on the composite score.}
\done{} Added an aggregation-robustness analysis: on Hubble-8B \dmech{} is rank~1 of~11 under all five
combinations (product, additive, harmonic mean, minimum, invariance-as-constraint). The 18-fold figure is
reframed as a convenience summary; the two effect sizes are the primary evidence.

\issue{6. Paraphrase evaluation too narrow for ``invariant''.}
\done{}\,/\,\pend{} Added a full second-benchmark replication (OpenUnlearning TOFU, whole-question
rephrasing, multi-token answers, at 1B/3B/8B), a different and harder paraphrase axis. A multi-paraphrase
sweep on Hubble itself remains future work and is stated as such.

\issue{7. Independence assumption questionable.}
\done{} Methodology and discussion now state that facts share checkpoints, templates, and overlapping
leave-one-out anchors, so the inferential statistics are \emph{benchmark-conditional}, not estimates over
an unrestricted population.

\issue{8. TOST insufficiently documented.}
\done{} TOST is de-emphasized: the equivalence bound is stated as exploratory and non-pre-registered, no
claim rests on it, and the Cliff's-delta effect sizes are made primary.

\issue{9. NPO ``erases the concept'' vs ``no method erases''.}
\done{} The internal contradiction is fixed: the semantic-residual result now reads ``NPO drives the
semantic residual lowest, toward the never-knew level'' (text and figure caption), not ``erases''.

\issue{10. Relearning: $\rho=-1.0$ at $n=4$.}
\done{} Reworded to ``descriptive corroboration, not independent validation'' with the entanglement at
$n=4$ stated. The confusing word ``shallow'' (which clashed with ``deep suppression'') is removed.

\issue{11. Tuned lens described inconsistently.}
\done{} The tuned lens is now explicitly an \emph{ordinal cross-check}. The passage that used its absolute
ranks is reworded to rely only on the ordinal agreement with the logit lens (Spearman $1.000$); we state we
do not read its absolute ranks as the true depth.

\issue{12. ``Causal Jacobian lens'' needs qualification.}
\done{} Renamed throughout to a \emph{first-order, interventionally-grounded readout}, a local
approximation around reference states, explicitly not a proof of causal necessity.

\issue{13. Hyperparameter detail.}
\done{} Appendix~A now has a hyperparameter table (Table~\ref{tab:hp} in the thesis): per-method learning
rate, epochs, retain weight, and method-specific coefficients (NPO $\beta$, RMU $\alpha$, steering
coefficient, intervention layers), plus the shared configuration and the forget-first selection criterion.

\issue{14. Possible tuning-on-evaluation.}
\done{} A direct ``Data reuse'' paragraph in the methodology now states that the same $44$ forget facts are
used for unlearning, hyperparameter selection, and every evaluation, and frames the study as an analysis of
the trained benchmark models rather than an out-of-sample estimate.

\issue{15. ``Retain'' pool naming.}
\done{} Methodology now states that ``Retain'' here denotes the never-knew reference pool, distinct from the
\emph{Keep}/retain-utility neighbour facts used to check collateral damage.

\issue{16. Reproducibility identifiers.}
\done{} The reproducibility section now lists the software environment (Python 3.11, PyTorch 2.11 / CUDA
12.8, Transformers 5.14, seed 0, deterministic decoding), the exact Hubble, unlearned, and TOFU checkpoint
identifiers, the Hubble revision, and the data sources. The code base is described as archived and
available on request rather than publicly released, per the reviewer's guidance for a private repository.

\section{The examiner questions}

\subsection*{On the core concept}
\begin{enumerate}[label=\arabic*.]
\item \textbf{Operational definition of erasure?} A fact counts as internally erased only if its
dual-anchor trajectory is indistinguishable from the never-knew reference across the band; no method reaches
this, so we report residual accessibility, observationally.
\item \textbf{Is linear decodability sufficient for knowledge?} No, and we do not claim it; the evidence is
comparative (same readout on knows and never-knew anchors) with tuned/Jacobian corroboration. Causal
patching is future work.
\item \textbf{Good rank from type/format priors?} The never-knew anchor shares those priors but lacks the
fact and sits far away ($\dmech=-1.36$); \dmech{} is the difference relative to both anchors, cancelling
shared priors.
\item \textbf{Why is never-knew a valid reference?} It is the strongest ground-truth lower bound (trained
identically minus the fact); used as a reference with leave-one-fact-out medians, not a perfect oracle.
\item \textbf{What if never-knew infers from context?} That raises the floor and shrinks separation, making
\dmech{} conservative; empirically the anchors separate cleanly.
\end{enumerate}

\subsection*{On \dmech{}}
\begin{enumerate}[label=\arabic*.,start=6]
\item \textbf{Why rank, not probability/log-odds/margin?} Rank is scale-free across five orders of
magnitude and robust to per-model calibration differences; probability fails invariance ($0.999$). We
still compare against probability and four MIA scores.
\item \textbf{Why absolute distance in log-rank?} The logarithm turns rank into a ratio; absolute distance
to each anchor, then subtraction, yields the ``closer to knows or never-knew'' reading that cancels the
common-mode paraphrase shift.
\item \textbf{Why average the band equally?} For simplicity and to avoid tuning; band-averaging is robust
(see band sensitivity below).
\item \textbf{Why layers 18--36, and pre-registered?} Fixed a priori as the top half where the logit lens
is reliable, not tuned on the outcome.
\item \textbf{Does the band choice change it?} Faithfulness under bands $12$--$30$, $18$--$36$, $24$--$36$
is $0.276$, $0.441$, $0.276$: the default is strongest, but the metric stays faithful (all far above the
best lexical $0.025$) and no band inverts the conclusion.
\item \textbf{Why median trajectories as anchors?} A robust central template; leave-one-fact-out avoids
self-comparison.
\item \textbf{Stability to anchor composition?} Leave-one-fact-out already tests fact-drop robustness; a
full anchor-subset sweep is a prepared backup.
\item \textbf{Positive \dmech{} because both anchors are close?} No: the never-knew anchor is far
($-1.36$) and the wrong-target pool is also negative. Precisely, positive \dmech{} means the measured
trajectory is closer to the knows reference than to the never-knew reference within the selected band; it is
an indicator of residual accessibility, not independent proof of stored knowledge. The single-anchor
min-rank fails, showing the two-anchor structure is doing the work.
\item \textbf{Decision threshold?} $\dmech>0$ (closer to knows) is the natural zero; we report
distributions and bootstrap CIs rather than a hard threshold, and judge faithfulness by effect sizes.
\end{enumerate}

\subsection*{On the statistical evaluation}
\begin{enumerate}[label=\arabic*.,start=15]
\item \textbf{Why Cliff's delta?} Non-parametric and scale-free, appropriate for bounded, skewed metrics;
the Wilcoxon $p$-value saturates at $n=44$.
\item \textbf{Why call $0.466$ invariant?} We no longer do (see Issue~4).
\item \textbf{TOST bounds?} Exploratory, non-pre-registered; de-emphasized.
\item \textbf{Independence with overlapping anchors?} Acknowledged; inference framed as
benchmark-conditional.
\item \textbf{Robustness across paraphrases?} TOFU uses whole-question rephrasing at three scales; a
Hubble multi-paraphrase sweep is future work.
\item \textbf{Why multiplicative?} One convenience choice; the ranking is invariant to product, additive,
harmonic, minimum, and constraint aggregations.
\item \textbf{Would the headline survive another formula?} Yes: rank~1 on Hubble-8B under all five, with
non-overlapping bootstrap CIs ($0.442\,[0.31,0.56]$ vs best lexical $0.025\,[0.00,0.07]$).
\end{enumerate}

\subsection*{On interpretation}
\begin{enumerate}[label=\arabic*.,start=22]
\item \textbf{Why does IDK count as retaining?} Because in a paired comparison its dual-anchor \dmech{}
exceeds Standard's on $100\%$ of facts, and it lies on the knows side of zero for $77\%$ (Issue~3),
independent of the coincident final rank.
\item \textbf{NPO ``erases'' vs no method erases?} The contradiction is fixed (Issue~9).
\item \textbf{Why does the deepest suppressor recover fastest?} A descriptive observation at $n=4$: RMU's
large-magnitude mid-stack perturbation is reversed by a few relearning steps.
\item \textbf{Tuned lens: knowledge or reconstruction?} We treat it as reconstruction-prone and use it only
as an ordinal cross-check.
\item \textbf{What makes the Jacobian causal?} It is interventionally grounded (the model's own
derivatives), a first-order local approximation, not a causal-necessity proof.
\item \textbf{Why does the Jacobian disagree with the logit lens for NPO?} The Jacobian keeps NPO buried
where the logit lens still sees a lingering linear trace: it separates a correlational trace from a
sensitivity-based signal. Reported as a difference between readouts, not an adjudicated truth.
\item \textbf{Which lens when they disagree?} The metric is built on the reproducible, parameter-free logit
lens; the other two are cross-checks and are all reported.
\end{enumerate}

\subsection*{On generalizability}
\begin{enumerate}[label=\arabic*.,start=29]
\item \textbf{Why only 44 birthdate facts?} That is Hubble's controlled birthdate forget set; TOFU adds
$80$ author facts at three scales.
\item \textbf{Multi-token or distributed concepts?} TOFU answers are multi-token sentences and \dmech{}
stays faithful; the nationality semantic-residual probe covers a distributed attribute.
\item \textbf{Real vs synthetic facts?} Hubble and TOFU are synthetic by design to provide ground-truth
anchors; real-data generalization is future work.
\item \textbf{Without knows/never-knew reference models?} The current metric needs both anchors; relaxing
that requirement is future work.
\item \textbf{Scale to production LLMs?} The paired-anchor requirement is the main barrier; noted as future
work.
\end{enumerate}

\section{The reviewer's ``must fix'' checklist}
\begin{enumerate}[label=\arabic*.]
\item Title page and supervisor. \done{}
\item Erasure to residual internal accessibility. \done{}
\item Do not call $0.466$ fully invariant. \done{}
\item Stop ``NPO erases the concept''. \done{}
\item Explain IDK vs Standard rank equality. \done{} (new dual-anchor distribution)
\item 18-fold result as a secondary summary. \done{}
\item State the sample size and single-paraphrase design early. \done{} (plus TOFU)
\item Tuned lens confirms ordering, not layer-wise computation. \done{}
\item Jacobian as first-order, interventionally-grounded. \done{}
\item Hyperparameters and data splits ready. \done{} (HP table + data-reuse paragraph in the thesis)
\end{enumerate}

\section{New evidence produced this session}
\begin{itemize}
\item \textbf{Aggregation robustness:} \dmech{} rank~1 on Hubble-8B under all five combinations; best
mechanistic metric and above every standard output metric on all five experiments.
\item \textbf{Second benchmark (TOFU 1B/3B/8B):} \dmech{} faithful at every scale (invar.\ shift
$0.35$--$0.42$, discrimination $0.91$--$0.95$), the best internal metric and above all standard output
metrics; a few robust lexical aggregates are competitive there, stated honestly.
\item \textbf{Bootstrap 95\% CIs:} Hubble-8B \dmech{} $[0.31,0.56]$ does not overlap the best lexical
$[0.00,0.07]$; TOFU-8B overlaps the truth ratio (statistically comparable), both far above probability.
These are fact-resampling intervals conditional on the fitted anchors, not a full anchor-uncertainty
bootstrap.
\item \textbf{Layer-band sensitivity:} faithful under layers $12$ to $30$, $18$ to $36$, $24$ to $36$; no
band inverts the result.
\item \textbf{IDK dual-anchor distribution:} median $+0.594$, on the knows side of zero for $77\%$ of facts
and, in a paired comparison, closer to knows than Standard on $100\%$ ($\delta=+0.997$,
$p=1.1\times10^{-13}$).
\end{itemize}

\section{Second-round corrections}
A later pass raised a further list of wording and documentation issues. All ten must-fix items are applied,
and the strongly-recommended items are addressed except one optional table row.

\subsection*{Must fix (all done)}
\begin{enumerate}[label=\arabic*.]
\item \textbf{``Causal readout''.} \done{} All remaining instances (background, methodology $\times2$) now
read \emph{first-order, interventionally-grounded readout}.
\item \textbf{Table ``more latent'' caption.} \done{} The tuned-lens table and figure captions now state
that we read only the ordering agreement (Spearman $1.000$), not the calibrated lens's absolute ranks.
\item \textbf{Future work vs completed TOFU.} \done{} Future work now records that the TOFU replication at
1B/3B/8B is done and points to further attributes and model families as the remaining extensions.
\item \textbf{``Single benchmark''.} \done{} Reworded to ``the primary Hubble analysis uses a single
attribute (birthdates)''.
\item \textbf{``Only faithful metric''.} \done{} Section title softened to ``Most Faithful Metric on
Hubble-8B''; a metric is now \emph{operationally} defined as faithful when discrimination $>0.5$ and
invariance shift $<0.5$, and on Hubble-8B only \dmech{} meets both.
\item \textbf{``IDK intact end to end''.} \done{} Softened to ``remains recoverable through to the output''.
\item \textbf{``Rules out the readout artifact''.} \done{} Softened to ``makes the early-layer-artifact
explanation unlikely''.
\item \textbf{``Output cannot separate methods''.} \done{} Reframed to ``does not separate'' / ``on these
output metrics the four methods look alike''.
\item \textbf{$0.2$ threshold.} \done{} Now stated as a heuristic noise margin set to $0.2$ for a $10$-way
infill at chance $0.1$.
\item \textbf{``Provably knows''.} \done{} Replaced with ``a model trained to memorise it (the knows
reference) and a model trained without it (the never-knew reference)'', calibrating against controlled
knowledge states.
\end{enumerate}

\subsection*{Strongly recommended}
\begin{itemize}
\item \textbf{Complete search grids + 1B settings.} \done{}\,/\,\pend{} The appendix now gives the full
grid (NPO lr $\in\{2\mathrm{e}{-5},5\mathrm{e}{-5}\}$, $\gamma\in\{0.5,1,2\}$; GradDiff lr
$\in\{5\mathrm{e}{-6},1\mathrm{e}{-5}\}$, $\gamma\in\{1,2,4\}$; IDK lr $\in\{2\mathrm{e}{-5},5\mathrm{e}{-5}\}$;
RMU steering $\in\{200,500,1000\}$, $\alpha\in\{50,100\}$; epochs checkpointed every two up to eight) and the
selected 8B values (Table~\ref{tab:hp} in the thesis). The 1B sweep uses the same grid and differs mainly in
the RMU intervention layers ($[3,4,5]$); the per-scale winners are stored with the artifacts.
\item \textbf{Paired paraphrase-change statistics.} \done{} Median Full$\to$Paraphrased change $-0.24$,
\dmech{} positive on every fact, paired Wilcoxon $p<0.001$ but small next to the anchor gap.
\item \textbf{TOFU multi-token teacher-forcing.} \done{} We teacher-force the preceding gold tokens and read
each answer token's rank at its own position; the trajectory is the mean rank over answer tokens, then
band-averaged.
\item \textbf{Separate discrimination.} \done{} On Hubble-8B $|\delta|=1.000$ against the never-knew pool
and $0.653$ against the wrong-target pool (combined $0.827$), so the gap is carried mainly by the never-knew
pool.
\item \textbf{Jacobian / tuned reference data.} \done{} The averaged Jacobian uses $24$ reference contexts
sampled from a pool of $38$ diverse generic sentences (leave-the-query-out); the tuned lens trains one
affine map per layer for $3000$ steps at learning rate $3\times10^{-4}$ (per-layer KL $\approx0.04$--$0.06$).
\item \textbf{Reduce ``18 times'' repetition.} \done{} Kept only in the abstract and the main results
statement; removed from the discussion, conclusion, and a figure caption.
\end{itemize}

\end{document}
